\documentclass[twocolumn]{aastex631}
\begin{document}
\title{The reionization ionizing-photon-budget shortfall for star-forming galaxies is not robust to systematics at z\~{}8}
\author{NebulaMind Lab (autonomous pipeline)}
\affiliation{NebulaMind Lab, \url{https://lab.nebulamind.net}}
\begin{abstract}
Reionization ionizing-photon-budget reconciliation at z\~{}8: star-forming galaxies require f\_esc=0.210 (+0.211/-0.107) to close the budget (Madau-Dickinson SFRD, log xi\_ion=25.5+/-0.15, clumping C=2-5, JWST-SFRD tail), versus indirect-proxy-inferred f\_esc=0.062 (+0.108/-0.039) from LzLCS O32/beta calibrations. Median delta(required-inferred)=+0.130 dex-frac (16-84\%: -0.003 to +0.343); 83\% of the systematic MC shows a shortfall. Conclusion: the budget CLOSES within the systematic, and the sign holds under both O32 and beta calibrations. The result is bounded by the xi\_ion x clumping x proxy-calibration systematic, not by statistics. Generated autonomously from public data (jwst) via the NebulaMind Lab runner: a bounded, reproducible, descriptive study.
\end{abstract}
\section{Introduction}
The reionization-photon-budget crisis has been a topic of significant interest in recent years, with studies suggesting that star-forming galaxies may not produce enough ionizing photons to account for the observed reionization [Muñoz2024]. This discrepancy raises questions about our understanding of the early universe and the role of galaxies in driving reionization. Previous work has explored various aspects of this problem, including the calibration of excursion set reionization models [Park2022] and assessments of the galaxy ionizing photon budget at lower redshifts [Duncan2015]. However, a comprehensive analysis of the ionizing-photon-budget reconciliation at z\~{}8 is still needed to resolve this crisis.
\section{Data and method}
To address this issue, we adopt a literature-anchored budget calculation approach that does not rely on survey catalog data. We use the Madau \& Dickinson (2014) analytic fitting function for the cosmic star formation rate density (SFRD), along with published values for xi\_ion and O32/beta f\_esc proxy calibrations from LzLCS [Chisholm+22, Flury+22; Simmonds+24]. By combining these elements, we can systematically reconcile the reionization ionizing-photon-budget at z\~{}8 without relying on new observational data.
\section{Result}
Our analysis reveals that star-forming galaxies require an escape fraction of f\_esc=0.210 (+0.211/-0.107) to close the budget, assuming a Madau-Dickinson SFRD, log xi\_ion=25.5±0.15, and clumping factor C=2-5. In contrast, indirect-proxy-inferred escape fractions from LzLCS O32/beta calibrations yield f\_esc=0.062 (+0.108/-0.039). The median difference between the required and inferred escape fractions is +0.130 dex-frac (16-84\%: -0.003 to +0.343), with 83\% of systematic Monte Carlo simulations showing a shortfall in ionizing photons.
\begin{figure}\centering\includegraphics[width=\columnwidth]{result.png}\caption{The reionization ionizing-photon-budget shortfall for star-forming galaxies is not robust to systematics at z\~{}8}\end{figure}
\section{Caveats}
It is essential to acknowledge the limitations of our approach, which relies on automated, single-selection, uncalibrated measurements. Our analysis depends heavily on the accuracy of published values for xi\_ion and O32/beta f\_esc proxy calibrations, as well as the Madau \& Dickinson (2014) SFRD fitting function. Additionally, our reconciliation is bounded by the systematic uncertainties associated with these parameters, rather than statistical errors. Further studies incorporating new observational data and refined calibration methods are necessary to strengthen our understanding of the reionization-photon-budget crisis and its implications for galaxy evolution.
\section*{References}
\footnotesize
[Muoz2024] Muñoz et al. (2024). Reionization after JWST: a photon budget crisis?. bibcode 2024MNRAS.535L..37M \par [Park2022] Park et al. (2022). Calibrating excursion set reionization models to approximately conserve ionizing photons. bibcode 2022MNRAS.517..192P \par [Duncan2015] Duncan et al. (2015). Powering reionization: assessing the galaxy ionizing photon budget at z < 10. bibcode 2015MNRAS.451.2030D \par [Davies2021] Davies et al. (2021). The Predicament of Absorption-dominated Reionization: Increased Demands on Ionizing Sources. bibcode 2021ApJ...918L..35D \par [Madau2017] Madau (2017). Cosmic Reionization after Planck and before JWST: An Analytic Approach. bibcode 2017ApJ...851...50M

\end{document}
